\Needspace{9\baselineskip}
\originalchapter{量子纠错条件与编码}
\originalsection{综合征怎样保护叠加}
直接测未知$a\ket0+b\ket1$会毁掉振幅. 编码把它映成$a\ket{0_L}+b\ket{1_L}$, 测的应是“发生哪种错误”的综合征, 不能区分逻辑0与1. 三位重复码$0_L=000,1_L=111$, 测$Z_1Z_2,Z_2Z_3$得$X_1,X_2,X_3$分别综合征$(-,+),(-,-),(+,-)$. 两逻辑码字对同一错误给相同综合征, 校正后振幅保留. 该码不能校正$Z$错误, 因$Z_1$在码内直接改变逻辑相位而综合征全正.

\begin{theorem}[Knill--Laflamme条件]\label{thm:KL}
码投影$P$可以完全纠正错误集合$\{E_a\}$, 当且仅当$PE_a^\dagger E_bP=c_{ab}P$, 其中$C=(c_{ab})$为与逻辑态无关的正半定矩阵.
\end{theorem}
\begin{proof}充分性: 对$C$作幺正对角化并同样重组错误,$F_a$, 得$PF_a^\dagger F_bP=d_a\delta_{ab}P$. $d_a>0$时$F_aP/\sqrt{d_a}$在码上等距且像空间互相正交. 测这些像空间、施其逆等距便恢复原逻辑态, 综合征概率$d_a$不含逻辑信息. 必要性: 若一个恢复过程把任意码态及其相干叠加都还原, 环境最终记录不能依赖逻辑态, 否则叠加相干会损失. 对联合等距写$\ket\psi\mapsto\sum_aE_a\ket\psi\ket a$, 环境矩阵元为$\bra\psi E_b^\dagger E_a\ket\psi$. 对所有码态独立, 再用极化恒等式把二次型恢复成矩阵元, 得题式.
\end{proof}

错误线性空间中的任意叠加也满足同条件. 因Pauli $I,X,Y,Z$张成单比特算符空间, 校正一个位置的这些离散错误便校正该位置的任意小环境耦合, 并非自然噪声必须真的是随机Pauli才可纠正. 退相干免费子空间则要求每个错误在码内为同一标量作用, 可无需主动检测; 它依赖特殊相关噪声, 不是所有局域噪声都可被动消除.

\originalsection{CSS与Steane码}
二进制线性码$C_2\subset C_1\subset\mathbb F_2^n$, 量子码字$\ket{x+C_2}=|C_2|^{-1/2}\sum_{y\in C_2}\ket{x+y}$,$x$取$C_1/C_2$陪集代表. $Z$类校验来自$C_1^\perp$, $X$类稳定子来自$C_2$; $u\cdot v=0$保证对易. 码编码$k=\dim C_1-\dim C_2$比特. Bit错误用$C_1$校验, Hadamard交换$X,Z$后phase错误用$C_2^\perp$校验; 不能误称直接用$C_2$的经典距离纠正所有phase.

Steane码取$C_1=[7,4,3]$Hamming码,$C_2=C_1^\perp=[7,3,4]$. 三个X及三个Z校验独立, 留$2^{7-6}=2$维码. 任何单个$X$或$Z$的综合征是校验矩阵对应非零列, 可定位错误; 任意单个$Y$同时有两类综合征. 故距离3可纠正一个任意位置错误. 两个不同位置的错误一般超出保证范围, 不能因分别会纠正X,Z就断言任意两个物理错误也可纠正.
